% Section 6.8, from lectures/L06/README.md: the objectives, the questions to test yourself, and the
% next lecture.

\needspace{10\baselineskip}
\section{Review}
\label{c6:sec:review}

\subsection*{What you should be able to do now}
After this chapter, you should be able to:
\begin{itemize}
  \item Build a 4-bit counter as a gate network, and explain, pointing at your own drawing, why
    nothing in it makes the count return to zero.
  \item Implement the same counter in VHDL, explain wraparound as a consequence of the register's
    bit width, and say why a simulator objects to a wrap the hardware performs happily.
  \item Implement a shift register in both SIPO and PISO configurations, know which fits which
    job, and know which direction this book shifts in and why sticking to one matters.
  \item Combine a counter with a shift register to detect that a complete word has arrived, and
    verify a clocked design by running its testbench.
\end{itemize}

\needspace{6\baselineskip}
\subsection*{Questions to test yourself}
You should be able to explain:
\begin{enumerate}
  \item Why does a \code{natural range 0 to 15} counter return to \code{0} after \code{15} with no
    explicit ``if at maximum, clear it'' logic? What has to be true of the range? Which element
    performs the wrap?
  \item A VHDL simulator aborts on the wrap your CircuitVerse drawing performs without complaint.
    Which of the two is modelling the hardware, and what does that tell you about trusting either
    alone?
  \item What is the practical difference between a SIPO and a PISO shift register? Which would you
    reach for to receive a stream of serial bits, and which to drive a chain of LEDs from a fixed
    pattern?
  \item A shift register and a counter are both ``N flip-flops sharing a clock''. What differs is
    what feeds each flip-flop's input. Describe that difference for each.
  \item In \code{serial_rx8}, why is \code{data_ready} a one-cycle pulse rather than a level that
    stays high once a byte has arrived?
  \item Why does \code{shift_enable} have to gate the bit counter as well as the shift itself?
  \item The board demonstration drives \code{shift_enable} from a button rather than slowing the
    clock down. Why is slowing the clock the wrong answer, and what would break if the button were
    wired to \code{shift_enable} without passing through \code{button_sync} first?
  \item \code{data_ready} needs a flip-flop of its own before an LED can show it. Which earlier
    chapter's problem is that, in different clothes?
\end{enumerate}

\needspace{6\baselineskip}
\subsection*{Looking ahead}
Chapter~\ref{c7:ch} is about timers:
\begin{itemize}
  \item Timers: a counter with a target count, which is what turns clock ticks into real time.
  \item Built by adding exactly two things to the counter you drew here, a comparator and a clear,
    so keep that CircuitVerse project.
  \item The walking LED: this chapter's shift register, paced by that timer and started by a
    synchronized button, built live in the course and run on the DE0-CV board.
\end{itemize}
